\documentclass{exam}
\usepackage{../../mypackages}
\usepackage{../../macros}

\title{Projets scientifiques - Empreinte carbone}
\author{N. Bancel}
\date{19 Novembre 2025}

% Mise en forme des solutions en bleu
\SolutionEmphasis{\color{blue}}
\renewcommand{\solutiontitle}{\noindent}

\begin{document}

\textbf{Collège Lycée Suger}
\hfill
\textbf{Projets scientifiques} \\

\textbf{Année 2025-2026}
\hfill
\textbf{Groupe 6ème-5ème} \par

{\let\newpage\relax\maketitle}

\begin{center}
  \textbf{Liste des choses qu'il reste à faire} \\
  Se répartir en différents groupes où chaque groupe est en charge d'une tâche spécifique
\end{center}

\section*{Cantine}

\subsection*{Interview du directeur de la cantine scolaire : Monsieur Foureau}


\begin{enumerate}[itemsep=2pt, topsep=2pt, parsep=0pt]
    \item \textbf{Quantités commandées :} Disposez-vous de données sur les quantités totales de nourriture commandées, idéalement en kilogrammes par type d'aliment ?
    \item \textbf{Matériel :} Quelle quantité de matériel est achetée chaque année (et de quel type ? Exemples : machines, vaisselle, etc)
    \item \textbf{Origine des produits :} D'où proviennent les produits alimentaires utilisés (local, régional, national, international) ?
    \item \textbf{Composition des menus :} Quelle est la proportion de repas contenant de la viande, du poisson ou des alternatives végétariennes ?
    \item \textbf{Composition des menus :} Avez-vous les menus à disposition pour les prochains mois / l'année 2026 ?
    \item \textbf{Labels durables :} Quel pourcentage des produits est issu de l'agriculture biologique ?
    \item \textbf{Saisonnalité :} Les fruits et légumes servis sont-ils majoritairement de saison ?
    \item \textbf{Ajustement des quantités :} Comment ajustez-vous les quantités préparées en fonction du nombre d’élèves présents ?
    \item \textbf{Déchets alimentaires :} Quelle quantité de déchets alimentaires est produite chaque jour ou chaque semaine (en kilogrammes) ? Existent-t-il des moyens de réduire cette quantité ?
    \item \textbf{Tri des déchets :} Quels types de déchets sont triés (plastique, carton, biodéchets, huiles, métal) et comment sont-ils valorisés ?
    \item \textbf{Emballages :} Les produits arrivent-ils en vrac ou emballés individuellement ?
    \item \textbf{Vaisselle jetable :} Utilisez-vous de la vaisselle jetable (gobelets, barquettes, couverts) ? À quelle fréquence ?
\end{enumerate}

\subsection*{Analyse du menu de la cantine}

\textbf{Analyse automatisée du menu :} Pouvez-vous créer, à l'aide d'un outil d'IA (ChatGPT, DeepSeek, MistralAI, Anthropic), un \emph{prompt} permettant d'analyser un menu de la cantine à partir d'un fichier PDF et de le reformater en un tableau indiquant :

\begin{itemize}[itemsep=1pt, topsep=2pt, parsep=0pt]
    \item la date,
    \item le moment de la journée,
    \item le type de plat (entrée, plat, dessert),
    \item la composition du plat,
    \item la catégorie du plat (viande, poisson, fruit ou légume, laitage, sucrerie, autre),
    \item si le plat est végétarien ou non.
\end{itemize}

\section*{Déplacements}

Un groupe est en charge de créer une version totalement finale du questionnaire pour les élèves en s'inspirant des 3 questionnaires réalisés. Il faudra ensuite envoyer un email à M. Barbarin et Mme Page pour les notifier que le questionnaire peut à nouveau être rempli par les élèves et leur demander de faire une relance.

\section*{Données sur l'établissement}

Créer une liste des choses restantes à faire pour l'inventaire dans les salles et notifier les élèves qui n'ont pas encore fait leur inventaire. Le document Google Sheet doit être prêt pour la semaine prochaine. 

\end{document}
